\documentclass[10pt,a4paper]{amsart}
\usepackage[margin=19mm]{geometry}
\usepackage{amsmath,amssymb,mathtools}
\usepackage[scr=rsfs,cal=euler]{mathalfa}
\usepackage{xfrac}
\usepackage[unicode,hidelinks,bookmarks=false]{hyperref}
\newcommand{\ie}{i.e.\ }
\newcommand{\cf}{cf.\ }
\newcommand{\resp}{respectively\ }
\newcommand{\Subsection}{\subsection}
\providecommand{\nobreakdash}{\nobreak\hbox{-}}
\setlength{\emergencystretch}{2em}
\multlinegap0pt
\usepackage{amssymb}
\usepackage[shortlabels]{enumitem}
\usepackage{dsfont}
\newtheorem{proposition}{Proposition}
\newtheorem{theorem}{Theorem}
\newtheorem{lemma}{Lemma}
\def\Fq#1{\mathbb{F}_q^{#1}}
\def\knorm#1#2{\lVert#2\rVert_{#1}}
\def\pnorm#1#2{|#2|_{#1}}
\title{The Mattila-Sj\"{o}lin problem for the $k$-distance over a finite field}
\author{Daewoong Cheong, Hunseok Kang,, Jinbeom Kim}
\date{Source: arXiv:2601.00529v1 (CC BY 4.0)}
\begin{document}
\maketitle

\noindent\emph{Source and licence.} The Mattila-Sj\"{o}lin problem for the $k$-distance over a finite field. Version arXiv:2601.00529v1 (CC BY 4.0), distributed by arXiv under the Creative Commons Attribution 4.0 International licence, http://creativecommons.org/licenses/by/4.0/, as recorded in the arXiv licence metadata for this exact version and shown on the abstract page https://arxiv.org/abs/2601.00529v1. Full licence notice: \texttt{licenses/arxiv-2601.00529-CC-BY-4.0.txt}.\par\smallskip

\noindent\textit{Editorial scope.} Counted unit: the article's theorem first* (source0.tex lines 425--430). The article's complete original proof of it is delivered with it (source0.tex lines 431--475, one \texttt{proof} environment of 45 lines). No mathematics is written, repaired or re-derived: the statement and the proof are the article's own lines, in the article's own notation, and no retained line is substituted or re-flowed.  Retained uncounted as the source-local context the proof needs: lines 390--397.  Nothing was omitted from any retained span, so the package carries no omission record.  No retained line contains a citation, so the wrapper carries no bibliography and no bibliography entry.  No retained line contains a figure environment, an image-inclusion command or drawing code, so no image is delivered and the package declares no asset dependency.  Editorial formatting only, in addition: an AMS \texttt{amsart} preamble that restates the article's own declarations and macros used by the retained lines, namely the environments \texttt{proposition}, \texttt{theorem}, \texttt{lemma} on separate counters, together with the standard packages those lines need.  The retained environments print in this document as Proposition 1, Theorem 1, Lemma 1 in order of appearance, whereas the article prints the same items with the numbering of its own full document.  The complete native source of the article as distributed in the arXiv e-print for this exact version is bundled unchanged as \texttt{sources/192006/source0.tex}.
\begin{proposition}\label{FSF} Let $t\in \mathbb F_q$. Then, for any $m\in \mathbb F_q^d,$ the Fourier transform $\widehat{S_{k}^t} (m)$ of the sphere can be written
$$\widehat{S_{k}^t} (m)=q^{-d-1} \bigg(A(m,t)+B(m)\bigg),$$
where 
\begin{align*}
    A(m,t) &=\sum_{\alpha = 0}^{k-1}\sum_{\substack{{I \subset [d]}\\{:|I| = d-\alpha}}}\sum_{s\in\Fq{*}}\chi(-st)\prod_{i \in I}\bigg(\eta(s)\mathcal{G}_1\chi\left(-\frac{{m}_i^2}{4s}\right)-1\bigg), \\
    B(m) &= \sum_{\alpha = 0}^{k-1}\sum_{\substack{{I \subset [d]}\\{:|I| = d-\alpha}}}(q-1)^{\mathcal{Z}(m_{I})}(-1)^{d-\alpha-\mathcal{Z}(m_I)}.
\end{align*}
\end{proposition}

\begin{theorem} \label{first*}
If $E \subset \mathbb{F}_q^d$ with $d \ge 2$, and 
$
|E| \ge C  q^{\max\left\{ \frac{d+1}{2},\ d - k \right\}},
$ then $D_k(E) = \mathbb{F}_q$.
\end{theorem}

\begin{proof}
Let \( d_E : E \times E \to \mathbb{F}_q \) denote the $k$-distance on $E$ defined by $
d_E(x,y) = \|x - y\|_{k}$ for $x,y \in E$.

Note that \( d_E \) is surjective if and only if for each \( t \in \mathbb{F}_q \), the set \( d_E^{-1}(t) \) is nonempty, or equivalently, the collection \( \{d_E^{-1}(t)\}_{t \in \mathbb{F}_q} \) forms a partition of \( E \times E \).

For $E \subset \Fq{d}$ and $t \in \Fq{}$, let $\nu_{E}(t)$ be the cardinality of the set $d_E^{-1}(t)$, i.e., $\nu_E(t)=\left|d_E^{-1} (t)\right|$. Then, to prove Theorem \ref{first*}, it suffices to show that $\nu_E(t) > 0$ for all  $t \in \Fq{}$. It is obvious that $\nu_E(0) > 0$ since $E$ is nonempty. Now it remains to  show  $\nu_E(t) > 0$
 for all nonzero $t \in \Fq{*}$.
By definition, $\nu_E(t)$ can be written
\begin{equation}\label{Nu}
\nu_{E}(t) = \sum_{\substack{{x,y\in E}\\{:\knorm{k}{x-y}=t}}}1 = \sum_{x,y\in E}S_{k}^t(x-y).
\end{equation}
Then we use  the Fourier inversion formula and the definition of the Fourier transform to obtain
\begin{align}\label{Nu1}
    \nu_E(t) & = \sum_{x,y \in \Fq{d}}E(x)E(y)\bigg{(}\sum_{m \in \Fq{d}}\chi(m\cdot(x-y))\widehat{S_{k}^t}(m)\bigg{)} \notag\\
    & = \sum_{m \in \Fq{d}} \widehat{S_{k}^t}(m) \sum_{x,y \in \Fq{d}}\overline{\chi(-m\cdot x)E(x)}\chi(-m\cdot y)E(y)\\
    & = q^{2d}\sum_{m\in\Fq{d}}\widehat{S_{k}^t}(m)|\widehat{E}(m)|^2.\notag
\end{align}
Invoking Proposition \ref{FSF},  we see that
 \begin{align}\label{AB}
\nu_{E}(t) &=q^{d-1}\sum_{m\in \Fq{d}}\pnorm{}{\widehat{E}(m)}^2\bigg(A(m,t)+B(m)
\bigg).
\end{align}
%where $A(m, t)$ and $B(m)$ were defined in Proposition \ref{FSF}.

Now we  require the following lemma, a key ingredient in deriving our main result,   contains  technically challenging terms that do not appear in the original MS problem.

\begin{lemma}\label{main} Let $A(m,t)$, and $B(m)$ be defined as in Proposition \ref{FSF}.  Then for every $E \subset \Fq{d}$, the following inequalities hold.
\begin{enumerate}
[label=\textup{(\roman*)}, ref=\textup{(\roman*)}]
\item \label{ABC1} $\left|\displaystyle{\sum_{m\in \Fq{d}}\pnorm{}{\widehat{E}(m)}^2}A(m,t)\right| \lesssim q^{-\frac{d-1}{2}}|E|,$
\item \label{ABC2} $\displaystyle{\sum_{m\in \Fq{d}}\pnorm{}{
\widehat{E}(m)}^2}B(m) \gtrsim q^{-d}|E|^2 - q^{-k}|E|$.
\end{enumerate}
\end{lemma}

With Lemma \ref{main} in hand, let us complete the proof of Theorem \ref{first*}. A proof of Lemma \ref{main} will be given in the next section.
The combination of \eqref{AB} and Lemma \ref{main} yields 
\begin{equation*}
	\nu_E(t) \gtrsim q^{-1}|E|^2 -q^{d-k-1}|E|-q^{\frac{d-1}{2}}|E|.
\end{equation*}
Observe that if \( |E| \gtrsim q^{d-k} \), then the term \( q^{-1}|E|^2 \) dominates \( q^{d-k-1}|E| \); similarly, if \( |E| \gtrsim q^{(d+1)/2} \), then \( q^{-1}|E|^2 \) dominates \( q^{(d-1)/2}|E| \). Thus, we obtain the desired conclusion that 
if $|E| \ge C  q^{\max\left\{ \frac{d+1}{2},\ d - k \right\}},$ then  $\nu_E(t) >0.$
This completes the proof.
\end{proof}

\end{document}
